\documentclass[11pt]{article}
% ===========================================================================
%  Shared preamble for the AP Calculus AB exam-preparation worksheets.
%  Input right after \documentclass. Files beginning with "_" are not compiled
%  on their own. Navy and gold, to match the course's lesson pages.
% ===========================================================================
\usepackage[letterpaper,top=2.3cm,bottom=2.2cm,left=1.7cm,right=1.7cm,headheight=28pt,footskip=30pt]{geometry}
\usepackage{amsmath,amssymb}
\usepackage[T1]{fontenc}
\usepackage[utf8]{inputenc}
\usepackage{lmodern}
\usepackage{xcolor}
\usepackage[most]{tcolorbox}
\usepackage{enumitem}
\usepackage{fancyhdr}
\usepackage{array}
\usepackage{tikz}
\usepackage{pgfplots}
\pgfplotsset{compat=1.18}

\definecolor{navy}{HTML}{161E54}
\definecolor{ink}{HTML}{27357E}
\definecolor{gold}{HTML}{B98A00}
\definecolor{goldbg}{HTML}{FFFDF3}
\definecolor{inkbg}{HTML}{F6F7FD}
\definecolor{rule}{HTML}{C9CFF0}
\definecolor{okgreen}{HTML}{1B7A44}
\definecolor{okbg}{HTML}{F1F8F3}

\pagestyle{fancy}
\fancyhf{}
\fancyhead[L]{\footnotesize Name: \underline{\hspace{4cm}}}
\fancyhead[C]{\textbf{\color{navy}Integration Academy}\\[-2pt]{\scriptsize Calculus AB \textperiodcentered{} AP\textsuperscript{\textregistered} Exam Preparation}}
\fancyfoot[L]{\tiny\color{black!55}\parbox[t]{0.78\textwidth}{AP\textsuperscript{\textregistered} is a registered trademark of the College Board, which was not involved in the production of, and does not endorse, this material. \textcopyright{} Integration Academy.}}
\fancyfoot[R]{\scriptsize\thepage}
\renewcommand{\headrulewidth}{1.2pt}
\renewcommand{\headrule}{\hbox to\headwidth{\color{gold}\leaders\hrule height \headrulewidth\hfill}}
\renewcommand{\footrulewidth}{0pt}
\setlength{\parindent}{0pt}
\setlength{\parskip}{4pt}

% Title block: \wstitle{code}{title}{unit line}
\newcommand{\wstitle}[3]{\begin{center}{\small\bfseries\color{gold}WORKSHEET #1}\\[3pt]{\LARGE\bfseries\color{navy}#2}\\[3pt]{\small\color{black!60}#3}\end{center}\vspace{2pt}}
% Section heading
\newcommand{\sech}[1]{\par\vspace{10pt}{\large\bfseries\color{navy}#1}\par\nobreak\vspace{2pt}{\color{rule}\hrule height0.8pt}\vspace{7pt}}
% A level tag printed beside a group of questions
\newcommand{\level}[1]{\par\vspace{6pt}{\footnotesize\bfseries\color{gold}\MakeUppercase{#1}}\par\nobreak\vspace{2pt}}

\newtcolorbox{ideabox}{colframe=ink,colback=inkbg,coltitle=white,colbacktitle=ink,fonttitle=\bfseries\small,title={Key ideas},arc=2pt,boxrule=0.8pt,left=8pt,right=8pt}
\newtcolorbox{example}[1]{colframe=rule,colback=white,coltitle=navy,colbacktitle=inkbg,fonttitle=\bfseries\small,title={#1},arc=2pt,boxrule=0.8pt,borderline west={3pt}{0pt}{ink},left=10pt,right=8pt}
\newtcolorbox{frbox}[1]{colframe=gold,colback=goldbg,coltitle=white,colbacktitle=gold,fonttitle=\bfseries\small,title={#1},arc=2pt,boxrule=0.8pt,breakable,left=8pt,right=8pt}
\newtcolorbox{solbox}{colframe=okgreen,colback=okbg,coltitle=white,colbacktitle=okgreen,fonttitle=\bfseries\small,title={Worked solutions},arc=2pt,boxrule=0.8pt,breakable,left=8pt,right=8pt}
\newtcolorbox{scorebox}[1]{colframe=navy,colback=white,coltitle=white,colbacktitle=navy,fonttitle=\bfseries\small,title={#1},arc=2pt,boxrule=0.8pt,breakable,left=8pt,right=8pt}

\newcommand{\soln}{\par\smallskip\textbf{\color{ink}Solution.}\ }
\newcommand{\result}[1]{\par\smallskip\textbf{#1}}
% blank working space of a given height
\newcommand{\work}[1]{\par\vspace{#1}}
% a point value in the margin of a scoring guide
\newcommand{\pts}[1]{\hfill{\footnotesize\bfseries\color{gold}[#1]}}

% Coordinate grid for graph questions: \begin{apgraph}{xmin}{xmax}{ymin}{ymax} ... \end{apgraph}
\newenvironment{apgraph}[4]{\begin{center}\begin{tikzpicture}\begin{axis}[
  axis lines=middle, xmin=#1, xmax=#2, ymin=#3, ymax=#4,
  grid=both, grid style={black!12}, xtick distance=1, ytick distance=1,
  ticklabel style={font=\scriptsize}, xlabel={$x$}, ylabel={$y$},
  width=10cm, height=7.2cm, samples=80, clip=false,
  every axis plot/.append style={line width=1.1pt, color=ink}]}{\end{axis}\end{tikzpicture}\end{center}}
\newcommand{\closeddot}[2]{\addplot[only marks, mark=*, mark size=2.2pt, color=ink] coordinates {(#1,#2)};}
\newcommand{\opendot}[2]{\addplot[only marks, mark=*, mark size=2.2pt, color=ink, mark options={fill=white}] coordinates {(#1,#2)};}

% Multiple-choice options: four across, or stacked when the options are long.
\newcommand{\opts}[4]{\par\vspace{3pt}\begin{tabular}{@{}p{0.235\linewidth}p{0.235\linewidth}p{0.235\linewidth}p{0.235\linewidth}@{}}(A)\ #1 & (B)\ #2 & (C)\ #3 & (D)\ #4\end{tabular}\par\vspace{9pt}}
\newcommand{\optsv}[4]{\par\vspace{3pt}\begin{tabular}{@{}l@{\ \ }p{0.9\linewidth}@{}}(A) & #1\\(B) & #2\\(C) & #3\\(D) & #4\end{tabular}\par\vspace{9pt}}

\fancyhead[R]{\footnotesize\color{black!60}Worksheet 3.5}
\begin{document}

\wstitle{3.5}{Choosing the Right Rule}{Unit 3: Composite, Implicit and Inverse Functions}

\begin{ideabox}
\begin{itemize}[leftmargin=*,itemsep=2pt,topsep=2pt]
  \item Ask: what is the \textbf{last operation} performed when evaluating the expression? That decides the first rule.
  \item A product of two functions: product rule. One function inside another: chain rule. Often both, one inside the other.
  \item A quotient can be rewritten as a product with a negative power when that is easier.
  \item Simplify before differentiating when the algebra is cheap; do not simplify afterwards unless asked.
\end{itemize}
\end{ideabox}

\sech{Worked examples}
\begin{example}{Example 1 \textperiodcentered{} Product outside, chain inside}
Find $\dfrac{d}{dx}\left[x^2e^{3x}\right]$.\soln $2xe^{3x}+x^2\cdot3e^{3x}$.\result{$e^{3x}\left(2x+3x^2\right)$.}
\end{example}

\begin{example}{Example 2 \textperiodcentered{} Chain outside, product inside}
Find $\dfrac{d}{dx}[\sin(x\cos x)]$.\soln $\cos(x\cos x)$ times the derivative of $x\cos x$.\result{$\cos(x\cos x)\,(\cos x-x\sin x)$.}
\end{example}

\begin{example}{Example 3 \textperiodcentered{} Quotient with a chain}
Find $\dfrac{d}{dx}\left[\dfrac{\ln\!\left(x^2+1\right)}{x}\right]$.\soln $\dfrac{x\cdot\frac{2x}{x^2+1}-\ln\!\left(x^2+1\right)}{x^2}$.\result{$\dfrac{2}{x^2+1}-\dfrac{\ln\!\left(x^2+1\right)}{x^2}$.}
\end{example}

\newpage
\sech{Practice}

\level{Core}
\begin{enumerate}[leftmargin=*,itemsep=2pt]
  \item Which rule is applied first to differentiate $y=(3x+1)^4\sin x$?\work{2.3cm}
  \item Which rule is applied first to differentiate $y=\sin\!\left((3x+1)^4\right)$?\work{2.3cm}
  \item Find $\dfrac{d}{dx}\left[e^{x}\ln x\right]$.\work{2.3cm}
  \item Find $\dfrac{d}{dx}\left[(1+\cos x)^2\right]$.\work{2.3cm}
\end{enumerate}

\level{Exam level}
\begin{enumerate}[leftmargin=*,itemsep=2pt,start=5]
  \item Find $\dfrac{d}{dx}\left[\sqrt{x}\,e^{-x}\right]$.\work{2.8cm}
  \item Find $\dfrac{d}{dx}\left[\tan^2(3x)\right]$.\work{2.8cm}
  \item Find $\dfrac{d}{dx}\left[\dfrac{x}{\sqrt{x^2+1}}\right]$ in simplified form.\work{2.8cm}
  \item Find $\dfrac{d}{dx}[\ln(\sec x+\tan x)]$ in simplified form.\work{2.8cm}
\end{enumerate}

\begin{frbox}{Question 9 \textperiodcentered{} Free-response practice}
Let $f(x)=x\sqrt{8-x^2}$ for $-2\sqrt2\le x\le2\sqrt2$.
\begin{enumerate}[label=(\alph*),leftmargin=*,itemsep=2pt]
  \item Show that $f'(x)=\dfrac{8-2x^2}{\sqrt{8-x^2}}$.\work{2.6cm}
  \item Find the values of $x$ at which the graph of $f$ has a horizontal tangent.\work{2.6cm}
  \item Write an equation of the tangent line at the point found in part (b) with $x>0$.\work{2.6cm}
\end{enumerate}
\end{frbox}

\newpage
\begin{solbox}
\begin{enumerate}[leftmargin=*,itemsep=4pt]
  \item The product rule: the last operation is the multiplication.
  \item The chain rule: the last operation is the sine.
  \item $e^{x}\ln x+\dfrac{e^{x}}{x}$.
  \item $2(1+\cos x)(-\sin x)=-2\sin x\,(1+\cos x)$.
  \item $\dfrac{e^{-x}}{2\sqrt x}-\sqrt x\,e^{-x}$.
  \item $2\tan(3x)\cdot\sec^2(3x)\cdot3=6\tan(3x)\sec^2(3x)$.
  \item $\dfrac{\sqrt{x^2+1}-\frac{x^2}{\sqrt{x^2+1}}}{x^2+1}=\dfrac{1}{\left(x^2+1\right)^{3/2}}$.
  \item $\dfrac{\sec x\tan x+\sec^2x}{\sec x+\tan x}=\sec x$.
  \item \begin{enumerate}[label=(\alph*),leftmargin=*,itemsep=3pt]
    \item $\sqrt{8-x^2}+x\cdot\dfrac{-x}{\sqrt{8-x^2}}=\dfrac{\left(8-x^2\right)-x^2}{\sqrt{8-x^2}}$.
    \item $8-2x^2=0$ gives $x=2$ and $x=-2$.
    \item $f(2)=2\sqrt4=4$ and the slope is 0: the line $y=4$.
  \end{enumerate}
\end{enumerate}
\end{solbox}

\end{document}
